\section{Problem}

The computational work is divided into two problems involving numerical linear
algebra and the numerical solution of a steady-state heat conduction problem.

\subsection{Problem 1: Symmetric Positive Definite Matrix}

Let $A \in \mathbb{R}^{10\times10}$ and $b \in \mathbb{R}^{10}$ be generated
according to the procedure given in Code 1 of Appendix~\ref{app:code1}, using
the student's enrollment number as the random seed.

The following tasks are considered:

\begin{enumerate}
    \item Determine whether the matrix $A$ is symmetric positive definite.
          If so, compute its Cholesky factorization
          \[
              A = GG^{T},
          \]
          and present the Cholesky factor $G$.

    \item Compute the condition number $\kappa(A)$ and justify the result.

    \item Solve the linear system
          \[
              Ax=b.
          \]
\end{enumerate}

\subsection{Problem 2: One-Dimensional Steady-State Heat Conduction}

Consider the one-dimensional steady-state heat conduction problem in a unit-length bar,

\begin{equation}
    -k\frac{d^2T}{dx^2}=f,
    \qquad 0\leq x\leq1,
\end{equation}

where $T(x)$ denotes the temperature, $k>0$ is the thermal conductivity,
and $f$ is a constant source term.

The boundary conditions are determined according to the procedure given in Code 2
of Appendix~\ref{app:code2}. One of the following cases is selected:

\begin{align*}
    \text{Case 1:}\quad &
    \left.\frac{dT}{dx}\right|_{x=0}=0,
    \qquad T(1)=0,
    \\[4pt]
    \text{Case 2:}\quad &
    T(0)=0,
    \qquad T(1)=0,
    \\[4pt]
    \text{Case 3:}\quad &
    \left.\frac{dT}{dx}\right|_{x=0}=0,
    \qquad
    \left.\frac{dT}{dx}\right|_{x=1}=0.
\end{align*}

Using the finite difference formulation specified in the assignment, with
$11$ grid points, the problem is represented by the linear system

\begin{equation}
    MT=s,
\end{equation}

where $M$ is the matrix associated with the discretized differential operator,
$s$ is the source vector including the boundary-condition contributions, and
$T$ is the vector of unknown temperatures.

The objectives are:

\begin{enumerate}
    \item Generate $M$ and $s$ using the student's enrollment number and solve
          the resulting linear system.
    \item Present the numerical temperature distribution $T(x)$.
    \item Plot the temperature distribution and interpret the result from a
          physical perspective.
\end{enumerate}